% Abstract Eng - Tóm tắt bằng Tiếng Anh
\mychapter{ABSTRACT}

{\fsize{12}
\textbf{Abstract:} This thesis proposes a digital identity architecture and a trust governance framework for digital academic credentials in multi-institution higher education, using Vietnam National University, Hanoi as a case study. The problem does not originate from a shortage of tools for detecting forged credentials, but from an architectural property of the current model: the information needed for verification resides with the issuing institution rather than within the credential itself, so the cost of verification scales with the number of verification requests. As a consequence, verification becomes an exception rather than a default action.

Building on the Self-Sovereign Identity (SSI) model, the work adopts Decentralized Identifiers (DID), Verifiable Credentials (VC) and a digital identity wallet in accordance with open W3C standards, combined with a hybrid on-chain and off-chain storage model.

The central contribution addresses the layer that existing systems leave unresolved. Cryptography establishes technical trust but cannot answer whether the organisation signing a credential is authorised to issue that particular type of credential. This gap is not theoretical: two published security analyses of deployed blockchain credential systems show it to be an exploitable class of vulnerability, and in both cases the attacked system had valid digital signatures and correctly anchored hashes. The thesis therefore specifies the authority chain --- issuance scope expressed as a bitmask, delegation validity periods, constraints on multi-level delegation, and revocation rules --- as rules enforced by smart contracts, so that authority state becomes public data any party can query independently rather than a record controlled by a single operator. The verification procedure comprises seven steps, including a comparison of credential content against the hash anchored on chain at issuance --- a step absent from systems that track status only, which consequently fail to detect content modified after issuance.

A prototype was implemented on a public network running the Move virtual machine, comprising three contract modules, two independent implementations of the indexing layer, and a non-custodial wallet running entirely in the browser. Of 126 contract unit tests, 83 construct a violating operation and assert that it is rejected with the expected error code. A four-experiment load measurement totalling 905,160 requests found no functional failures and median latency one to two orders of magnitude below the stated threshold at moderate load; the same measurement also showed that tail latency cannot be measured reliably when the load generator shares a host with the system under test. The precise statement of the result is that the specified rules were not violated in the scenarios tested --- not a proof of correctness.

The research further derives three design principles whose scope extends beyond digital credentials: revocation must carry weaker preconditions than issuance; an irreversible operation must have its precondition proven by an action of the affected party rather than by a declaration of the acting party; and restoring a previously revoked role requires a versioning mechanism in order to preserve the audit trail.

Four limitations are stated explicitly. Separating on-chain from off-chain data \emph{reduces} exposure of personal data but does not \emph{guarantee} the right to erasure, since an unsalted hash on a public chain remains personal data under current regulatory interpretation. Each issuer signs with a single key, so a compromised key permits valid issuance across the issuer's entire authority scope. The tail-latency threshold requirement is not verifiable on the available measurement setup. And the thesis concludes on technical feasibility, not on usability --- two questions that do not follow from one another.

Within the scope of this thesis, {``}academic credential{''} refers to the set of a learner's academic records, including degrees, transcripts, certificates and other learning outcomes. {``}Digital academic credential{''} denotes the digital form of these records, modelled within the SSI architecture as a verifiable credential.

\vspace{0.5cm}
\textit{\textbf{Keywords:}} \textit{Trust Governance, Digital Credentials, Decentralized Identity, Self-Sovereign Identity, Trust Registry, Smart Contract.}
}
